% !TEX root = ../algebra_02.tex

\section{Построение свободной группы}

Пусть $n \in \mathbb{N}$.
Зададим алфавит $A = \{x_1, x_1', \dots, x_n, x_n'\}$, где символ $x_i'$ играет роль формального обратного элемента для $x_i$.
Определим множество слов $W$:
\[ W = \{a_1 \dots a_k \mid k \ge 0, a_i \in A, i = 1, \dots, k\} \]
Введем операцию конкатенации (умножения) слов $W \times W \to W$:
\[ (a_1 \dots a_k, b_1 \dots b_l) \mapsto a_1 \dots a_k b_1 \dots b_l \]

Будем говорить, что слово $\widetilde{w}$ получается из слова $w$ \textbf{вставкой}, если $w = w_0 w_1$ и $\widetilde{w} = w_0 x_i x_i' w_1$ или $\widetilde{w} = w_0 x_i' x_i w_1$ для некоторых $w_0, w_1 \in W$.
Слово $\widetilde{w}$ получается из $w$ \textbf{сокращением}, если $w$ получается из $\widetilde{w}$ вставкой.
Введем отношение $\sim$: $w \sim \widetilde{w}$, если $\widetilde{w}$ получается из $w$ конечным числом вставок или сокращений.
Это отношение является отношением эквивалентности.

Определим множество $F_n$ как фактормножество $W$ по отношению эквивалентности $\sim$:
\[ F_n := W / \sim \]
Обозначим через $[w]$ класс эквивалентности слова $w$.
Введем умножение на $F_n$:
\[ F_n \times F_n \to F_n, \quad ([w], [w']) \mapsto [ww'] \]

\begin{Lemma}{Корректность умножения}{}
	Операция умножения на $F_n$ задана корректно, то есть результат не зависит от выбора представителей классов эквивалентности.
\end{Lemma}

\begin{proof}
	Пусть $w \sim \widetilde{w}$ и $w' \sim \widetilde{w}'$. Тогда выполнена цепочка эквивалентностей:
	\[ ww' \sim \widetilde{w}w' \sim \widetilde{w}\widetilde{w}' \implies ww' \sim \widetilde{w}\widetilde{w}' \]
	Следовательно, классы совпадают: $[ww'] = [\widetilde{w}\widetilde{w}']$.
\end{proof}

\begin{Theorem}{Конструкция свободной группы}{}
	Множество $F_n = W / \sim$ является свободной группой с $n$ свободными образующими $f_1, \dots, f_n$, где $f_i := [x_i]$ для $i = 1, \dots, n$.
\end{Theorem}

\begin{proof}
	Сначала покажем, что $F_n$ образует группу:
	\begin{itemize}
		\item Ассоциативность операции очевидна из ассоциативности конкатенации слов.
		\item Нейтральным элементом является класс пустого слова $[]$.
		\item Проверим существование обратных элементов. Заметим, что:
		\[ [x_i] \cdot [x_i'] = [x_i x_i'] = [] \]
		\[ [x_i'] \cdot [x_i] = [x_i' x_i] = [] \]
		В общем случае, пусть $g \in F_n$. Представим его как $g = [a_1 \dots a_m] = [a_1] \dots [a_m]$, где $a_i \in A$. Тогда обратным элементом к $g$ будет класс $[a_m]^{-1} \dots [a_1]^{-1}$, где под $[x_i]^{-1}$ понимается $[x_i']$, а под $[x_i']^{-1}$ понимается $[x_i]$. В этом случае:
		\[ g \cdot \left([a_m]^{-1} \dots [a_1]^{-1}\right) = [] \quad \text{и} \quad \left([a_m]^{-1} \dots [a_1]^{-1}\right) \cdot g = [] \]
	\end{itemize}

	Теперь проверим универсальное свойство. Пусть задана группа $G$ и элементы $g_1, \dots, g_n \in G$. Определим отображение $\Phi: W \to G$ на образующих:
	\[ \Phi(x_i) = g_i, \quad \Phi(x_i') = g_i^{-1} \]
	И продолжим его на произвольные слова:
	\[ \Phi(a_1 \dots a_m) = \Phi(a_1) \dots \Phi(a_m) \]

	Покажем, что если $w \sim w'$, то $\Phi(w) = \Phi(w')$. Достаточно проверить это для одной операции вставки или сокращения. Пусть $\widetilde{w} = v_0 x_i x_i' v_1$, тогда:
	\[ \Phi(v_0 x_i x_i' v_1) = \Phi(v_0) g_i g_i^{-1} \Phi(v_1) = \Phi(v_0)\Phi(v_1) = \Phi(v_0 v_1) \]
	Следовательно, отображение корректно спускается на фактормножество, то есть существует отображение $\varphi: F_n \to G$ такое, что $\varphi([w]) = \Phi(w)$. Проверим, что $\varphi$ — гомоморфизм:
	\[ \varphi([w][w']) = \varphi([ww']) = \Phi(ww') = \Phi(w)\Phi(w') = \varphi([w])\varphi([w']) \]

	Докажем единственность. Пусть существует другой гомоморфизм $\varphi': F_n \to G$, удовлетворяющий условию $\varphi'(f_i) = g_i$. Тогда:
	\[ \varphi'([x_i]) = g_i = \varphi([x_i]) \]
	\[ \varphi'([x_i']) = g_i^{-1} = \varphi([x_i']) \]
	Так как значения на всех образующих и их обратных совпадают, то для любого $g \in F_n$ выполнено $\varphi'(g) = \varphi(g)$.
\end{proof}

\subsection*{Несократимые слова}

\begin{Definition}{Несократимое слово}{}
	Слово $w \in W$ называется несократимым, если в нем нельзя произвести ни одного сокращения (то есть оно не содержит подслов вида $x_i x_i'$ или $x_i' x_i$).
\end{Definition}

\begin{Proposition}{О несократимых словах}{}
	Каждый класс эквивалентности (элемент $F_n = W / \sim$) содержит ровно одно несократимое слово.
\end{Proposition}

\begin{proof}
	\textbf{Существование:} очевидно, что в каждом классе эквивалентности есть хотя бы одно несократимое слово, так как длина слова конечна, и процесс последовательных сокращений завершается за конечное число шагов.

	\textbf{Единственность:} пусть $w \sim w'$, где $w$ и $w'$ — несократимые слова. Предположим от противного, что $w \neq w'$. По определению отношения эквивалентности существует конечная последовательность преобразований:
	\[ \lambda: w = w_0 \mapsto w_1 \mapsto \dots \mapsto w_n = w' \]
	где каждый переход $w_i \mapsto w_{i+1}$ является либо вставкой, либо сокращением подслова $x_i x_i^{-1}$. Определим длину последовательности преобразований $\lambda$ как сумму длин всех участвующих в ней слов:
	\[ l(\lambda) = l(w_0) + l(w_1) + \dots + l(w_n) \]
	Среди всех возможных последовательностей преобразований, переводящих $w$ в $w'$, выберем ту, для которой значение $l(\lambda)$ минимально. Далее в выбранной последовательности выберем слово $w_j$ максимальной длины. Если таких слов несколько, выберем слово с наименьшим индексом $j \ge 1$.

	Поскольку начальное слово $w_0 = w$ несократимо, первый шаг $w_0 \mapsto w_1$ обязательно является вставкой, следовательно, локальный максимум длины существует и $j \ge 1$. Рассмотрим переходы в окрестности максимума: $w_{j-1} \mapsto w_j$ (обязательно вставка) и $w_j \mapsto w_{j+1}$ (обязательно сокращение).

	Из комбинаторики слов следует, что можно произвести оба сокращения из $w_j$ (те, которые вели бы обратно в $w_{j-1}$ и в $w_{j+1}$) независимым образом. Это означает, что существует слово $w_j'$, к которому можно прийти из $w_{j-1}$ сокращением, и из которого можно прийти в $w_{j+1}$ вставкой. При этом $l(w_j') = l(w_j) - 4$.

	Тогда мы можем заменить фрагмент пути $w_{j-1} \mapsto w_j \mapsto w_{j+1}$ на $w_{j-1} \mapsto w_j' \mapsto w_{j+1}$, что строго уменьшит общую сумму длин слов $l(\lambda)$. Это противоречит выбору последовательности $\lambda$ с минимальным значением $l(\lambda)$. Следовательно, $w = w'$.
\end{proof}

\begin{Remark}{}{}
	Ввиду доказанного утверждения, можно считать, что свободная группа $F_n$ состоит в точности из несократимых слов в алфавите $A$. При этом образующий $x_i$ отождествляется с $f_i$, а $x_i'$ — с $f_i^{-1}$.
\end{Remark}
